\documentclass[12pt,a4paper]{article}
\usepackage[T2A]{fontenc}
\usepackage[utf8]{inputenc}
\usepackage[russian]{babel}
\usepackage{amsmath}
\usepackage{geometry}
\geometry{margin=2cm}
\setlength{\parindent}{0pt}
\setlength{\parskip}{0.45em}

\begin{document}

\begin{center}
{\Large Финансы и кредит}\\
{\large Вариант 1: краткий учебный разбор по локальным материалам}
\end{center}

\textbf{Важно.} Это не готовая работа ``под сдачу'', а короткий конспект, чтобы быстро решить вариант по тем же шаблонам, что были на семинарах.

\textbf{На что опирался из локальных файлов}
\begin{itemize}
\item \texttt{Фонд оценочных средств ФиК 2026.pdf}: задачи 1, 3, 4 и 5 почти совпадают по типу.
\item \texttt{ФиК Практика 2 весна 2026 СДО.pptx}: есть транспортный налог, авансовые платежи и правило учета месяцев регистрации.
\item \texttt{ФиК\_Практика\_Модуль\_3\_Кредит\_весна\_26.pdf}: есть овердрафт и аннуитет.
\item \texttt{ФиК\_Практика\_Международные\_финансы.pdf}: есть формулы по call-опциону и премии.
\item \texttt{finance\_credit\_kr\_prep.pdf}: короткий свод всех нужных формул.
\end{itemize}

\textbf{1. Дефицит бюджета}

Из семинаров: сначала находим доходы, потом сравниваем их с расходами.

\[
\text{Собственные доходы} = 0{,}85 \cdot 870
\]
\[
\text{Доходы всего} = 870 + 0{,}85 \cdot 870
\]
\[
\text{Дефицит} = 1680 - \text{Доходы всего}
\]

Если итог положительный, это дефицит. Если отрицательный, это уже профицит.

\textbf{2. Транспортный налог}

По практике 2:
\begin{itemize}
\item для легкового автомобиля мощностью 175 л.с. берется ставка 50 руб./л.с.;
\item месяц регистрации и месяц снятия с регистрации считаются полными;
\item авансовые платежи считаются за 1, 2 и 3 кварталы;
\item по семинарскому шаблону квартальный аванс равен $1/4$ годовой суммы налога с учетом коэффициента.
\end{itemize}

В этой задаче автомобиль числится с января по ноябрь включительно, то есть:
\[
K_{\text{год}}=\frac{11}{12}
\]

Так как в 1, 2 и 3 кварталах автомобиль был зарегистрирован весь квартал, то:
\[
\text{АП}_{1\text{кв}}=\frac{175 \cdot 50}{4}
\]
\[
\text{АП}_{2\text{кв}}=\frac{175 \cdot 50}{4}
\]
\[
\text{АП}_{3\text{кв}}=\frac{175 \cdot 50}{4}
\]
\[
\text{Налог за год}=175 \cdot 50 \cdot \frac{11}{12}
\]

Если преподаватель просит именно \textit{доплату по итогам года}, то из годовой суммы надо вычесть три авансовых платежа.

\textbf{3. Ипотека}

Дано: квартира 5 млн руб., первоначальный взнос 20\%, ставка 12\% годовых, срок 6 лет.

\[
S = 5\,000\,000 \cdot (1-0{,}20)=4\,000\,000
\]
\[
r=\frac{12\%}{12}=1\%=0{,}01
\]
\[
N=6 \cdot 12=72
\]
\[
A=\frac{S \cdot r}{1-(1+r)^{-N}}
\]

Где:
\begin{itemize}
\item $S$ --- сумма кредита;
\item $r$ --- месячная ставка;
\item $N$ --- число ежемесячных платежей;
\item $A$ --- ежемесячный аннуитетный платеж.
\end{itemize}

Дальше:
\[
\text{Общая сумма выплат}=A \cdot 72
\]
\[
\text{Сумма процентов за весь срок}=A \cdot 72 - 4\,000\,000
\]

\textbf{4. Овердрафт}

\textbf{Сначала проверь условие на фото.} Там строка про документы к оплате выглядит не совсем однозначно: при остатке на счете 150 млн руб. сумма сделки почти наверняка тоже должна быть в \textit{млн руб.}, а не в тыс. руб., иначе задача не бьется по смыслу.

Если брать семинарский шаблон и считать, что к оплате поступило \textbf{170 млн руб.}, то:
\[
\text{Сумма овердрафта}=170-150=20 \text{ млн руб.}
\]

Дальше удобно идти \textbf{строго по таблице из практики}:
\begin{enumerate}
\item Сначала считаем проценты до первого поступления:
\[
I_1=\frac{20 \cdot 16 \cdot 10}{100 \cdot 365}
\]
\item Из первого поступления 5 млн руб. сначала покрываются проценты $I_1$, остаток идет в погашение тела овердрафта.
\item Получаем новый остаток долга $O_1$.
\item Дальше по семинарскому образцу считаем:
\[
I_2=\frac{O_1 \cdot 16 \cdot 12}{100 \cdot 365}
\]
\item После второго поступления 10 млн руб. снова сначала закрываются проценты $I_2$, затем тело долга. Получаем новый остаток $O_2$.
\item Затем:
\[
I_3=\frac{O_2 \cdot 16 \cdot 15}{100 \cdot 365}
\]
\item Из третьего поступления 12 млн руб. сначала гасим $I_3$, остаток направляем на основной долг.
\end{enumerate}

\textbf{Замечание.} В локальной практике по задачам с несколькими поступлениями преподавательский шаблон считает проценты на текущий остаток долга и число дней до соответствующего поступления. Если захочешь переписать решение, лучше держаться именно этого образца.

\textbf{5. Проданный call-опцион}

Здесь важно не перепутать роли: в задаче \textbf{инвестор продал} call, то есть он --- \textbf{продавец опциона}.

Премия, которую он получает сразу:
\[
\text{Премия}=0{,}3 \cdot 40\,000
\]

Если к моменту исполнения:

\textit{Ситуация 1.} $S_T=86{,}5$ руб./долл., страйк $K=85{,}9$.

Для продавца call финансовый результат:
\[
\Pi=\bigl(0{,}3-(86{,}5-85{,}9)\bigr)\cdot 40\,000
\]

\textit{Ситуация 2.} $S_T=85{,}6$ руб./долл.

Так как курс ниже страйка, покупателю невыгодно исполнять call. Экономически корректный результат для продавца:
\[
\Pi=0{,}3 \cdot 40\,000
\]

Максимальная прибыль продавца call всегда ограничена полученной премией.

\textbf{Что проверить перед оформлением}
\begin{itemize}
\item В задаче 2 уточни, что именно просит преподаватель: налог за год или доплату по итогам года после авансов.
\item В задаче 4 перепроверь единицы измерения в строке про сумму сделки.
\item В задаче 5 безопаснее показать и логику из семинара, и короткую фразу, что невыгодный call покупатель не исполняет.
\end{itemize}

\end{document}
